\subsection{Ranking and Selection Metrics}
\label{app:routing-metrics}
Let $L_q=(s_1,\ldots,s_K)$ be the ordered skill list for task $q$, and let $G_q$ be its required skill set. Recall at cutoff $k$ measures the fraction of required skills in the first $k$ positions:
\[
\operatorname{Recall@}k(q)=\frac{|\{s_1,\ldots,s_k\}\cap G_q|}{|G_q|}.
\]
For nDCG, the relevance of each skill is binary. With $y_{q,j}=\mathds{1}[s_j\in G_q]$, we compute
\[
\operatorname{DCG@}k(q)=\sum_{j=1}^{k}\frac{y_{q,j}}{\log_2(j+1)},
\]
\[
\operatorname{nDCG@}k(q)=\frac{\operatorname{DCG@}k(q)}{\sum_{j=1}^{\min(k,|G_q|)}1/\log_2(j+1)}.
\]
Thus, nDCG@1 is the fraction of tasks whose first-ranked skill is required. Recall@1 can be lower when a task requires multiple skills, because it divides by the size of the required set.

We distinguish the ranked list from the selected set $P_q$. Set-level F1 is
\[
\operatorname{F1}(q)=\frac{2|P_q\cap G_q|}{|P_q|+|G_q|}.
\]
For strict predicate routing, $P_q$ contains only skills whose predicates evaluate to true. A top-$k$ routing list also contains predicate-negative skills when they fall within the ranked cutoff. The set-level F1 reported for predicate ablations therefore measures the selected set, rather than the first $k$ entries of the complete ranked list.

Dataset scores average the corresponding per-task metric. The overall score pools task instances across datasets; it is not an unweighted average of the six dataset scores. When comparing retrievers, we first compute this task-weighted score for each retriever and then average across BM25, Vector, and SkillSight.
